% Part: sets-functions-relations
% Chapter: sets
% Section: important-sets

\documentclass[../../../include/open-logic-section]{subfiles}

\begin{document}

\olfileid[tr]{sfr}{set}{imp}
\olsection{Bazı Önemli Kümeler}

\begin{ex}
Çoğunlukla elemanları matematiksel nesneler olan kümelerle
uğraşacağız. Bu kümelerden dördü, özel adları olacak kadar önemlidir:
\begin{multline*}
    \Nat = \{0, 1, 2, 3, \ldots\} \\
    \shoveright{\text{doğal sayılar kümesi}}\\
    \shoveleft{\Int = \{\ldots, -2, -1, 0, 1, 2, \ldots\}} \\
    \shoveright{\text{tam sayılar kümesi}}\\
    \shoveleft{\Rat = \Setabs{\nicefrac{m}{n}}{m, n \in \Int\text{ ve }n \neq 0}}\\
    \shoveright{\text{rasyonel sayılar kümesi}}\\
    \shoveleft{\Real = (-\infty, \infty)}\\
    \text{gerçel sayılar kümesi (kontinuum)}
\end{multline*}
Bunların hepsi \emph{sonsuz} kümelerdir; yani her birinin sonsuz sayıda
elemanı vardır.

Bu kümelerde ilerledikçe sayı dağarcığımıza \emph{daha çok} sayı ekleriz.
Gerçekten $\Nat \subseteq \Int \subseteq \Rat \subseteq \Real$ olduğu açık
olmalıdır: her doğal sayı bir tam sayıdır; her tam sayı bir rasyonel sayıdır;
her rasyonel sayı da bir gerçel sayıdır. Benzer biçimde
$\Nat \subsetneq \Int \subsetneq \Rat$ olduğu da açık olmalıdır; çünkü $-1$
bir tam sayı olup doğal sayı değildir ve $\nicefrac{1}{2}$ rasyonel olup tam
sayı değildir. $\Rat \subsetneq \Real$ olduğu, yani rasyonel olmayan bazı
gerçel sayıların bulunduğu daha az açıktır\oliflabeldef{sfr:arith:real:realline}{;
ancak buna \olref[arith][real]{realline} kesiminde döneceğiz}{}.

Bazen $\PosInt = \{1, 2, 3, \dots\}$ pozitif tam sayılar kümesini ve yalnızca
ilk iki doğal sayıyı içeren $\Bin = \{0, 1\}$ kümesini de kullanacağız.
\end{ex}

\begin{tagblock}{compsci}
\begin{ex}[Dizgiler]
Başka ilginç bir örnek, bir $A$ alfabesi üzerindeki \emph{sonlu dizgilerin}
$A^{*}$ kümesidir: $A$ alfabesinin elemanlarından oluşan her sonlu dizi, $A$ üzerinde
bir dizgidir. Her $A$ alfabesi için \emph{boş dizgiyi $\Lambda$} da $A$
üzerindeki dizgilere katarız. Örneğin
\begin{multline*}
\Bin^*
=\{\Lambda,0,1,00,01,10,11,\\
000,001,010,011,100,101,110,111,0000,\ldots\}.
\end{multline*}
$x=x_{1}\ldots x_{n}\in A^{*}$, $A$ alfabesinden alınmış $n$ ``harften'' oluşan bir
dizgiyse dizginin \emph{uzunluğunun} $n$ olduğunu söyler ve $\len{x}=n$
yazarız.
\end{ex}
\end{tagblock}

\begin{ex}[Sonsuz diziler]
Herhangi bir $A$ kümesi için, $A$ kümesinin elemanlarından oluşan
sonsuz dizilerin $A^\omega$ kümesini de ele alabiliriz. Bir
$a_1a_2a_3a_4\dots$ sonsuz dizisi, her biri $A$ kümesinin elemanı olan
nesnelerin tek yönde sonsuza uzanan bir listesidir.
\end{ex}

\end{document}
